% TOP_PROBLEM: {"id":30007692,"problem_number":"LOCAL-30007692","title":"Scientific discovery bottlenecks","category_id":21,"category":{"id":21,"name":"economics","display_name":"Economics","description":"Open economic theory statements, published research agendas, and proposed scoped specifications; question provenance and historical priority are qualified per record.","slug":"economics","order_index":21,"created_at":"2026-10-08T00:00:00Z"},"status":"research_question","external_url":null,"legacy_ids":[],"published":false,"statement":"Locate remaining pharmaceutical discovery bottlenecks by comparing AI assistance across hypothesis generation, testing and independent validation.","economics":{"original_id":181,"field":"Innovation and AI economics","kind":"empirical","formalization_basis":"adapted_research_question","source_status":"explicit_open","priority_status":"earliest_found","priority_note":"Earliest relevant explicit question or agenda passage located in this audit. No claim of first-ever formulation; earlier versions and later solutions were not exhaustively traced.","earliest_verified_reference":{"authors":["Erik Brynjolfsson","Anton Korinek","Ajay K. Agrawal"],"title":"A Research Agenda for the Economic Implications of Transformative Artificial Intelligence","year":2024,"url":"https://digitaleconomy.stanford.edu/app/uploads/2024/11/ETAI-White-Paper.pdf","doi":null,"locator":"Section 3.2, printed pp. 9–10, scientific discovery bottleneck questions","evidence_summary":"Explicitly asks which discovery stages remain bottlenecks as AI automates research."},"current_reference":{"authors":["Erik Brynjolfsson","Anton Korinek","Ajay K. Agrawal"],"title":"A Research Agenda for the Economic Implications of Transformative Artificial Intelligence","year":2024,"url":"https://digitaleconomy.stanford.edu/app/uploads/2024/11/ETAI-White-Paper.pdf","doi":null,"locator":"Section 3.2, printed pp. 9–10, scientific discovery bottleneck questions","evidence_summary":"Explicitly asks which discovery stages remain bottlenecks as AI automates research."},"formalization_note":"The source verifies a broader question or research agenda; the population, estimand, equations and restrictions here are our scoped adaptation. This is not a quoted canonical conjecture, and the audit does not prove contemporary unsolved status for every special case."},"statusQualification":"Published question or research agenda verified at the cited locator. The mathematical specification is adapted for this catalogue and includes additional assumptions; source publication does not establish first-ever priority or certify this exact model remains unsolved. The source verifies a broader question or research agenda; the population, estimand, equations and restrictions here are our scoped adaptation. This is not a quoted canonical conjecture, and the audit does not prove contemporary unsolved status for every special case.","statusReviewedAt":"2026-10-08"}
% FIRST_POSTED: 2026-10-08
% ENTRY_KIND: statement_only
% !TeX program = lualatex
\documentclass[11pt]{article}
\usepackage{fontspec}
\usepackage[a4paper,margin=25mm]{geometry}
\usepackage{amsmath,amssymb,mathtools}
\usepackage{xurl}
\usepackage[hidelinks]{hyperref}
\newcommand{\sref}[2]{\hyperlink{source:#1:#2}{[S#2]}}
\setlength{\emergencystretch}{3em}
\begin{document}
% problemId:30007692
\section*{Scientific discovery bottlenecks}\label{30007692}
\subsection{Definitions and mathematical statement}
Consider pharmaceutical discovery projects entering a prespecified laboratory network during one year. Each project passes through hypothesis generation \(H\), experimental testing \(E\), and independent validation \(V\), with staff hours \(c_s\) and processing capacity \(k_s\) at stage \(s\in\{H,E,V\}\). An intervention \(a=(a_H,a_E,a_V)\in\{0,1\}^3\) supplies a fixed AI assistance package at selected stages; zero denotes the existing workflow. Let \(N(a)\) be the number of discoveries meeting a blinded, preregistered reproducibility and clinical relevance threshold within twelve months, and let \(C(a)>0\) include labor, computation, experiments and validation costs. For the distribution of eligible projects and laboratories, define productivity \(P(a)=E[N(a)]/E[C(a)]\). A stage is a marginal bottleneck if a small feasible increase in its capacity produces the largest derivative of expected validated discoveries per additional dollar. Estimate the factorial intervention effects on \(P(a)\), stage waiting times and marginal capacity values, allowing faster hypothesis generation to congest testing or validation. Laboratory randomization, recorded queues and independent outcome adjudication are proposed inputs; within-laboratory interactions must be retained. Which stages remain binding, and when does investment in a downstream stage create larger validated gains than further AI improvement upstream? The target is this laboratory population, rather than all scientific discovery.

\par\textbf{Formulation and scope.} The source verifies a broader question or research agenda; the population, estimand, equations and restrictions here are our scoped adaptation. This is not a quoted canonical conjecture, and the audit does not prove contemporary unsolved status for every special case.

\par\textbf{Open-question provenance.} An explicit question or research agenda was verified at the source locator. This does not by itself certify that every later version remains unresolved \sref{30007692}{1}.

\par\textbf{Historical priority.} Earliest relevant explicit question or agenda passage located in this audit. No claim of first-ever formulation; earlier versions and later solutions were not exhaustively traced.

\subsection{Short English statement}
Locate remaining pharmaceutical discovery bottlenecks by comparing AI assistance across hypothesis generation, testing and independent validation.

\subsection{Sources}
\begin{itemize}\raggedright
\item[\textnormal{[S1]}]\hypertarget{source:30007692:1}{} Erik Brynjolfsson, Anton Korinek, Ajay K. Agrawal. 2024. A Research Agenda for the Economic Implications of Transformative Artificial Intelligence. Section 3.2, printed pp. 9–10, scientific discovery bottleneck questions.
\url{https://digitaleconomy.stanford.edu/app/uploads/2024/11/ETAI-White-Paper.pdf}.
{\small\textit{Earliest verified question/agenda reference; historical priority qualified below. Explicitly asks which discovery stages remain bottlenecks as AI automates research.}}
\end{itemize}
\end{document}
